% GENERATED FILE. Edit canonical lessons, then run npm run generate:books.
% canonical-source-hash: fnv1a64:14c05d2cc687621a
% canonical-lessons: TA-C398-pokku, TA-C398-caval, TA-C398-catanai, TA-C398-nilaittanmai, TA-C398-vivakaram

\chapter{Trends and Challenges}
\label{ch:ta-trends-and-challenges}
\hlchaptermodality{398}

\begin{chapteropening}
\textbf{By the end of this chapter:} \emph{I can talk about trends, challenges and achievements.}

Complete the last lesson of chapter 398: I can talk about trends, challenges and achievements.
\end{chapteropening}

\section[\texorpdfstring{\ta{போக்கு} (pōkku)}{போக்கு (pōkku)}]{\ta{போக்கு} (pōkku) — a trend; a tendency}

\label{lesson:TA-C398-pokku}



\begin{quote}
Before the new one: say the Tamil for personal, then the Tamil for rational thought.
\end{quote}

\subsection*{You'll want to know: \ta{போக்கு}}
\textbf{\ta{போக்கு}} — \emph{pōkku} — \textquotedblleft{}a trend; a tendency\textquotedblright{}.

\begin{quote}
\ta{கைப்பேசியில்} \ta{படிக்கும்} \ta{போக்கு} \ta{இப்போது} \ta{அதிகரித்து} \ta{வருகிறது}.
\end{quote}

\begin{quote}
\emph{kaippēciyil paṭikkum pōkku ippōtu atikarittu varukiṟatu.}
\end{quote}

\begin{quote}
The trend of reading on phones is growing now.
\end{quote}

Also a way of behaving: \ta{அவன்} \ta{போக்கு} \ta{சரியில்லை}, his ways are not right.

\begin{cousinweb}
\textbf{\ta{போக்கு}} (\emph{pōkku}) is native Dravidian: DEDR 4572 lists Tamil \emph{pō}, \textquotedblleft{}to go\textquotedblright{}, and \emph{pōkku}, \textquotedblleft{}going; to cause to go\textquotedblright{}: the way things are going.
\end{cousinweb}

\begin{grammarlens}[title={What you've built}]
One more word for saying what happens if.
\end{grammarlens}

\subsection*{Your turn}
\begin{itemize}
\raggedright
  \item \emph{Say it:} \emph{pōkku}
  \item \emph{Say it:} \emph{pōkku}, once more
  \item \emph{Say it:} \emph{kaippēciyil paṭikkum pōkku ippōtu atikarittu varukiṟatu.} — The trend of reading on phones is growing now.
  \item \emph{Recall:} say the Tamil for personal, then the Tamil for rational thought, then say \emph{pōkku} again
\end{itemize}

\begin{tcolorbox}[breakable,colback=teal!4,colframe=teal!35!black,title={Before you move on}]
What does \ta{போக்கு} mean? (\textquotedblleft{}A trend; a tendency\textquotedblright{}.) Say it once more.
\end{tcolorbox}
\section[\texorpdfstring{\ta{சவால்} (cavāl)}{சவால் (cavāl)}]{\ta{சவால்} (cavāl) — a challenge}

\label{lesson:TA-C398-caval}



\begin{quote}
Before the new one: say the Tamil for rational thought, then the Tamil for a trend.
\end{quote}

\subsection*{You'll want to know: \ta{சவால்}}
\textbf{\ta{சவால்}} — \emph{cavāl} — \textquotedblleft{}a challenge\textquotedblright{}.

\begin{quote}
\ta{தண்ணீர்} \ta{இல்லாவிட்டால்} \ta{விவசாயம்} \ta{ஒரு} \ta{பெரிய} \ta{சவால்}.
\end{quote}

\begin{quote}
\emph{taṇṇīr illāviṭṭāl vivacāyam oru periya cavāl.}
\end{quote}

\begin{quote}
Without water, farming is a big challenge.
\end{quote}

\ta{சவால்} \ta{விடு}, to throw down a challenge.

\begin{cousinweb}
\textbf{\ta{சவால்}} (\emph{cavāl}) is usually traced to Arabic \emph{suʾāl}, \textquotedblleft{}a question, a demand\textquotedblright{}, which reached Tamil through Hindustani, where it also means a challenge.
\end{cousinweb}

\begin{grammarlens}[title={What you've built}]
One more word for saying what happens if.
\end{grammarlens}

\subsection*{Your turn}
\begin{itemize}
\raggedright
  \item \emph{Say it:} \emph{cavāl}
  \item \emph{Say it:} \emph{cavāl}, once more
  \item \emph{Say it:} \emph{taṇṇīr illāviṭṭāl vivacāyam oru periya cavāl.} — Without water, farming is a big challenge.
  \item \emph{Recall:} say the Tamil for rational thought, then the Tamil for a trend, then say \emph{cavāl} again
\end{itemize}

\begin{tcolorbox}[breakable,colback=teal!4,colframe=teal!35!black,title={Before you move on}]
What does \ta{சவால்} mean? (\textquotedblleft{}A challenge\textquotedblright{}.) Say it once more.
\end{tcolorbox}
\section[\texorpdfstring{\ta{சாதனை} (cātaṉai)}{சாதனை (cātaṉai)}]{\ta{சாதனை} (cātaṉai) — an achievement; a record}

\label{lesson:TA-C398-catanai}



\begin{quote}
Before the new one: say the Tamil for a trend, then the Tamil for a challenge.
\end{quote}

\subsection*{You'll want to know: \ta{சாதனை}}
\textbf{\ta{சாதனை}} — \emph{cātaṉai} — \textquotedblleft{}an achievement; a record\textquotedblright{}.

\begin{quote}
\ta{இந்த} \ta{முறை} \ta{வென்றால்} \ta{அவள்} \ta{புதிய} \ta{சாதனை} \ta{படைப்பாள்}.
\end{quote}

\begin{quote}
\emph{inta muṟai veṉṟāl avaḷ putiya cātaṉai paṭaippāḷ.}
\end{quote}

\begin{quote}
If she wins this time, she will set a new record.
\end{quote}

\ta{சாதனை} \ta{படை}, to set a record; \ta{சாதனையாளர்}, an achiever.

\begin{cousinweb}
\textbf{\ta{சாதனை}} (\emph{cātaṉai}) is borrowed from Sanskrit \emph{sādhana}, \textquotedblleft{}furthering, accomplishing, completion\textquotedblright{} (CDIAL 13334), with the Tamil ending -\ta{ஐ}.
\end{cousinweb}

\begin{grammarlens}[title={What you've built}]
One more word for saying what happens if.
\end{grammarlens}

\subsection*{Your turn}
\begin{itemize}
\raggedright
  \item \emph{Say it:} \emph{cātaṉai}
  \item \emph{Say it:} \emph{cātaṉai}, once more
  \item \emph{Say it:} \emph{inta muṟai veṉṟāl avaḷ putiya cātaṉai paṭaippāḷ.} — If she wins this time, she will set a new record.
  \item \emph{Recall:} say the Tamil for a trend, then the Tamil for a challenge, then say \emph{cātaṉai} again
\end{itemize}

\begin{tcolorbox}[breakable,colback=teal!4,colframe=teal!35!black,title={Before you move on}]
What does \ta{சாதனை} mean? (\textquotedblleft{}An achievement; a record\textquotedblright{}.) Say it once more.
\end{tcolorbox}
\section[\texorpdfstring{\ta{நிலைத்தன்மை} (nilaittaṉmai)}{நிலைத்தன்மை (nilaittaṉmai)}]{\ta{நிலைத்தன்மை} (nilaittaṉmai) — stability}

\label{lesson:TA-C398-nilaittanmai}



\begin{quote}
Before the new one: say the Tamil for a challenge, then the Tamil for an achievement.
\end{quote}

\subsection*{You'll want to know: \ta{நிலைத்தன்மை}}
\textbf{\ta{நிலைத்தன்மை}} — \emph{nilaittaṉmai} — \textquotedblleft{}stability\textquotedblright{}.

\begin{quote}
\ta{அரசியலில்} \ta{நிலைத்தன்மை} \ta{இருந்தால்} \ta{நாடு} \ta{வளரும்}.
\end{quote}

\begin{quote}
\emph{araciyalil nilaittaṉmai iruntāl nāṭu vaḷarum.}
\end{quote}

\begin{quote}
If there is political stability, the country will grow.
\end{quote}

\ta{நிலை}, a firm standing + \ta{தன்மை}, nature.

\begin{cousinweb}
\textbf{\ta{நிலைத்தன்மை}} (\emph{nilaittaṉmai}) is native: \emph{nilai}, \textquotedblleft{}a standing, firmness\textquotedblright{} (DEDR 3675), and \emph{taṉmai}, \textquotedblleft{}nature, quality\textquotedblright{} (DEDR 3196): the quality of standing firm.
\end{cousinweb}

\begin{grammarlens}[title={What you've built}]
One more word for saying what happens if.
\end{grammarlens}

\subsection*{Your turn}
\begin{itemize}
\raggedright
  \item \emph{Say it:} \emph{nilaittaṉmai}
  \item \emph{Say it:} \emph{nilaittaṉmai}, once more
  \item \emph{Say it:} \emph{araciyalil nilaittaṉmai iruntāl nāṭu vaḷarum.} — If there is political stability, the country will grow.
  \item \emph{Recall:} say the Tamil for a challenge, then the Tamil for an achievement, then say \emph{nilaittaṉmai} again
\end{itemize}

\begin{tcolorbox}[breakable,colback=teal!4,colframe=teal!35!black,title={Before you move on}]
What does \ta{நிலைத்தன்மை} mean? (\textquotedblleft{}Stability\textquotedblright{}.) Say it once more.
\end{tcolorbox}
\section[\texorpdfstring{\ta{விவகாரம்} (vivakāram)}{விவகாரம் (vivakāram)}]{\ta{விவகாரம்} (vivakāram) — a matter; an affair}

\label{lesson:TA-C398-vivakaram}



\begin{quote}
Before the new one: say the Tamil for an achievement, then the Tamil for stability.
\end{quote}

\subsection*{You'll want to know: \ta{விவகாரம்}}
\textbf{\ta{விவகாரம்}} — \emph{vivakāram} — \textquotedblleft{}a matter; an affair\textquotedblright{}.

\begin{quote}
\ta{இந்த} \ta{விவகாரம்} \ta{நீதிமன்றத்துக்குப்} \ta{போனால்} \ta{நீண்ட} \ta{காலம்} \ta{ஆகும்}.
\end{quote}

\begin{quote}
\emph{inta vivakāram nītimaṉṟattukkup pōṉāl nīṇṭa kālam ākum.}
\end{quote}

\begin{quote}
If this matter goes to court, it will take a long time.
\end{quote}

Often a disputed matter; in speech also a fuss.

\begin{cousinweb}
\textbf{\ta{விவகாரம்}} (\emph{vivakāram}) is borrowed from Sanskrit \emph{vyavahāra}, \textquotedblleft{}behaviour, custom, business\textquotedblright{} (CDIAL 12174), whose Prakrit form also means \textquotedblleft{}lawsuit\textquotedblright{}. Tamil writes vya as vi and h as \ta{க}.
\end{cousinweb}

\begin{grammarlens}[title={What you've built}]
One more word for saying what happens if.
\end{grammarlens}

\subsection*{Your turn}
\begin{itemize}
\raggedright
  \item \emph{Say it:} \emph{vivakāram}
  \item \emph{Say it:} \emph{vivakāram}, once more
  \item \emph{Say it:} \emph{inta vivakāram nītimaṉṟattukkup pōṉāl nīṇṭa kālam ākum.} — If this matter goes to court, it will take a long time.
  \item \emph{Recall:} say the Tamil for an achievement, then the Tamil for stability, then say \emph{vivakāram} again
\end{itemize}

\begin{tcolorbox}[breakable,colback=teal!4,colframe=teal!35!black,title={Before you move on}]
What does \ta{விவகாரம்} mean? (\textquotedblleft{}A matter; an affair\textquotedblright{}.) Say it once more.
\end{tcolorbox}
